% =============================================================================
%  Comandos, ambientes e estilos usados no texto
% =============================================================================
\usepackage[most]{tcolorbox}

% ---- notação
\newcommand{\vetor}[1]{\bm{#1}}
\newcommand{\vr}{\vetor{r}}
\newcommand{\vv}{\vetor{v}}
\newcommand{\va}{\vetor{a}}
\newcommand{\vh}{\vetor{h}}
\newcommand{\ve}{\vetor{e}}
\newcommand{\dd}{\mathrm{d}}
\newcommand{\ii}{\mathrm{i}}
\newcommand{\UA}{\,\mathrm{UA}}
\newcommand{\dia}{\,\mathrm{dia}}
\newcommand{\km}{\,\mathrm{km}}
\newcommand{\sgn}{\operatorname{sgn}}
\newcommand{\codigo}[1]{\texttt{\small #1}}
\newcommand{\arquivo}[1]{\texttt{\small #1}}

% ---- teoremas (numeração por capítulo)
\theoremstyle{plain}
\newtheorem{proposicao}{Proposição}[section]
\theoremstyle{definition}
\newtheorem{definicao}{Definição}[section]
\newtheorem{exemplo}{Exemplo}[section]
\theoremstyle{remark}
\newtheorem*{observacao}{Observação}

% ---- caixas didáticas
% "Para começar": leitura de entrada, sem cálculo (ensino médio)
\newtcolorbox{paracomecar}[1][]{enhanced,breakable,colback=blue!3,colframe=blue!45!black,
  boxrule=0.5pt,arc=1mm,left=2mm,right=2mm,top=1mm,bottom=1mm,
  fonttitle=\bfseries\small,title={Para começar#1}}
% "Para aprofundar": detalhe técnico que pode ser pulado na primeira leitura
\newtcolorbox{paraaprofundar}[1][]{enhanced,breakable,colback=orange!4,colframe=orange!60!black,
  boxrule=0.5pt,arc=1mm,left=2mm,right=2mm,top=1mm,bottom=1mm,
  fonttitle=\bfseries\small,title={Para aprofundar#1}}
% "Lição de projeto": um erro real que o projeto cometeu e corrigiu
\newtcolorbox{licao}[1][]{enhanced,breakable,colback=red!3,colframe=red!50!black,
  boxrule=0.5pt,arc=1mm,left=2mm,right=2mm,top=1mm,bottom=1mm,
  fonttitle=\bfseries\small,title={Lição de projeto#1}}
% "No programa": onde a ideia aparece na tela
\newtcolorbox{noprograma}[1][]{enhanced,breakable,colback=green!3,colframe=green!40!black,
  boxrule=0.5pt,arc=1mm,left=2mm,right=2mm,top=1mm,bottom=1mm,
  fonttitle=\bfseries\small,title={No programa#1}}

% ---- código-fonte
\definecolor{fundocodigo}{RGB}{248,248,250}
\definecolor{comentario}{RGB}{90,110,90}
\definecolor{palavra}{RGB}{20,60,150}
\definecolor{textocod}{RGB}{150,40,40}
\lstdefinelanguage{TypeScript}{
  keywords={const,let,var,function,return,if,else,for,while,break,continue,new,
    export,import,from,type,interface,readonly,as,of,in,throw,true,false,null,undefined},
  morekeywords=[2]{number,string,boolean,void,Float64Array,Math},
  sensitive=true,
  morecomment=[l]{//},
  morecomment=[s]{/*}{*/},
  morestring=[b]',
  morestring=[b]",
  morestring=[b]`
}
\lstset{
  basicstyle=\ttfamily\scriptsize,
  backgroundcolor=\color{fundocodigo},
  keywordstyle=\color{palavra}\bfseries,
  keywordstyle=[2]\color{palavra},
  commentstyle=\color{comentario}\itshape,
  stringstyle=\color{textocod},
  numbers=left, numberstyle=\tiny\color{gray}, numbersep=6pt,
  frame=single, rulecolor=\color{gray!40}, framexleftmargin=4pt,
  breaklines=true, breakatwhitespace=false, columns=fullflexible, keepspaces=true,
  showstringspaces=false, tabsize=2, captionpos=t,
  literate={á}{{\'a}}1 {é}{{\'e}}1 {í}{{\'i}}1 {ó}{{\'o}}1 {ú}{{\'u}}1
           {Á}{{\'A}}1 {É}{{\'E}}1 {Í}{{\'I}}1 {Ó}{{\'O}}1 {Ú}{{\'U}}1
           {ã}{{\~a}}1 {õ}{{\~o}}1 {Ã}{{\~A}}1 {Õ}{{\~O}}1
           {â}{{\^a}}1 {ê}{{\^e}}1 {ô}{{\^o}}1 {Â}{{\^A}}1 {Ê}{{\^E}}1 {Ô}{{\^O}}1
           {ç}{{\c c}}1 {Ç}{{\c C}}1 {à}{{\`a}}1 {ü}{{\"u}}1
           {—}{{---}}1 {−}{{-}}1 {·}{{$\cdot$}}1 {×}{{$\times$}}1 {≤}{{$\le$}}1 {≥}{{$\ge$}}1
           {π}{{$\pi$}}1 {μ}{{$\mu$}}1 {√}{{$\surd$}}1 {°}{{$^\circ$}}1 {≈}{{$\approx$}}1
}
\renewcommand{\lstlistingname}{Código}
\renewcommand{\lstlistlistingname}{Lista de códigos}

% ---- algoritmos
\SetAlgoCaptionSeparator{ --}
\SetKwComment{Com}{$\triangleright$\ }{}
\SetAlFnt{\small}
% palavras-chave em português (independe da versão do algorithm2e)
\SetKwInput{KwIn}{Entrada}
\SetKwInput{KwOut}{Saída}
\SetKwFor{For}{para}{faça}{fim}
\SetKwFor{ForEach}{para cada}{faça}{fim}
\SetKwFor{While}{enquanto}{faça}{fim}
\SetKwIF{If}{ElseIf}{Else}{se}{então}{senão se}{senão}{fim}
\SetKw{KwTo}{até}
\SetKw{Return}{retorne}
